\chapter{Running jobs with Slurm}
\label{ch:slurm}

\section{Two ways to run}\label{two-ways-to-run}

{\def\LTcaptype{none} % do not increment counter
\begin{longtable}[]{@{}
  >{\raggedright\arraybackslash}p{(\linewidth - 4\tabcolsep) * \real{0.3333}}
  >{\raggedright\arraybackslash}p{(\linewidth - 4\tabcolsep) * \real{0.3333}}
  >{\raggedright\arraybackslash}p{(\linewidth - 4\tabcolsep) * \real{0.3333}}@{}}
\toprule\noalign{}
\begin{minipage}[b]{\linewidth}\raggedright
\end{minipage} & \begin{minipage}[b]{\linewidth}\raggedright
\texttt{srun\ ...\ -\/-pty\ bash}
\end{minipage} & \begin{minipage}[b]{\linewidth}\raggedright
\texttt{sbatch\ script.sh}
\end{minipage} \\
\midrule\noalign{}
\endhead
\bottomrule\noalign{}
\endlastfoot
Style & Interactive: you get a shell on a compute node & Hand it in and walk away \\
Good for & Short tests, installs & All real work \\
If you disconnect & Work stops & It carries on \\
\end{longtable}
}

\section{A first job}\label{a-first-job}

Save this as \texttt{first.sh}:

\begin{lstlisting}[language=bash]
#!/bin/bash
# The first line (the "shebang") says which program should run this file.

# Lines starting #SBATCH are instructions to Slurm. They must come before any command.
#SBATCH --job-name=first             # a label for the queue
#SBATCH --time=00:05:00              # limit: hours:minutes:seconds. The job is stopped at this point
#SBATCH --cpus-per-task=2            # cores for this job
#SBATCH --mem=2G                     # memory for this job
#SBATCH --output=first_%j.out        # where printed output goes; %j becomes the job ID

hostname                             # which node am I on?
nproc                                # how many cores did I get?
date                                 # when did it run?
\end{lstlisting}

\begin{lstlisting}[language=bash]
# cluster
sbatch first.sh                # hand it in. Prints: Submitted batch job <jobid>
squeue -u $USER                # my jobs. PD means pending, R means running
cat first_<jobid>.out          # read what it printed
\end{lstlisting}

When your job vanishes from \texttt{squeue}, it has ended, either successfully or not, so always read the log. By default a job starts in the folder from which you submitted it.

\section{The settings you need}\label{the-settings-you-need}

{\def\LTcaptype{none} % do not increment counter
\begin{longtable}[]{@{}
  >{\raggedright\arraybackslash}p{(\linewidth - 4\tabcolsep) * \real{0.3333}}
  >{\raggedright\arraybackslash}p{(\linewidth - 4\tabcolsep) * \real{0.3333}}
  >{\raggedright\arraybackslash}p{(\linewidth - 4\tabcolsep) * \real{0.3333}}@{}}
\toprule\noalign{}
\begin{minipage}[b]{\linewidth}\raggedright
Setting
\end{minipage} & \begin{minipage}[b]{\linewidth}\raggedright
Meaning
\end{minipage} & \begin{minipage}[b]{\linewidth}\raggedright
Choosing it
\end{minipage} \\
\midrule\noalign{}
\endhead
\bottomrule\noalign{}
\endlastfoot
\texttt{-\/-job-name} & A label & Short and useful \\
\texttt{-\/-time} & Hard limit, also \texttt{days-hh:mm:ss} & Run a test, and add a margin. An accurate limit helps a job start sooner \\
\texttt{-\/-cpus-per-task} & Cores for a threaded program & Only what it can actually use \\
\texttt{-\/-ntasks} & Separate processes, as in MPI & 1, unless you use MPI \\
\texttt{-\/-nodes} & Number of nodes & Usually leave it, unless using MPI \\
\texttt{-\/-mem} or \texttt{-\/-mem-per-cpu} & Memory & Check real use with \texttt{sacct}, then trim it \\
\texttt{-\/-partition} or \texttt{-p} & Which queue & See below \\
\texttt{-\/-output}, \texttt{-\/-error} & Files for printed output and errors & Use \texttt{\%x} job name, \texttt{\%j} job ID, \texttt{\%A} array ID, \texttt{\%a} task number. The folder must already exist, or the log is lost \\
\texttt{-\/-array} & Many tasks; see the next chapter & \\
\texttt{-\/-gres=gpu:\textless{}type\textgreater{}:1} & A GPU & Only on GPU partitions, and only if your code uses it \\
\texttt{-\/-mail-type}, \texttt{-\/-mail-user} & Email on start and end & Handy for long jobs; avoid it for arrays \\
\end{longtable}
}

The definitive list is in the \href{https://slurm.schedmd.com/sbatch.html}{\texttt{sbatch} manual}\footnote{\texttt{https://slurm.schedmd.com/sbatch.html}}.

Partitions, \textbf{Kelvin2 only}, from \href{https://ni-hpc.github.io/nihpc-documentation/Running\%20Jobs/}{Running Jobs}\footnote{\texttt{https://ni-hpc.github.io/nihpc-documentation/Running\%20Jobs/}}; check there for changes:

{\def\LTcaptype{none} % do not increment counter
\begin{longtable}[]{@{}lll@{}}
\toprule\noalign{}
Partition & Time limit & For \\
\midrule\noalign{}
\endhead
\bottomrule\noalign{}
\endlastfoot
\texttt{k2-hipri} (default) & 3 hours & Short jobs \\
\texttt{k2-medpri} & 24 hours & Longer jobs \\
\texttt{k2-lowpri} & Not listed & Long jobs \\
\texttt{k2-himem} & 3 days & Very large memory \\
\texttt{k2-gpu-*} & 3 days & GPUs \\
\texttt{k2-sandbox} & See \texttt{sinfo} & Testing, per the Green HPC page \\
\end{longtable}
}

\section{Watching and managing}\label{watching-and-managing}

\begin{lstlisting}[language=bash]
# cluster
squeue -u $USER                 # my queued and running jobs
squeue -u $USER -h | wc -l      # -h removes the header, so this counts my jobs
scancel <jobid>                 # cancel a job
scancel -u $USER                # cancel all my jobs. Careful
sinfo -s                        # a summary of partitions
sacct -j <jobid> --format=JobID,Elapsed,MaxRSS,State    # after it ends: time, peak memory, outcome
\end{lstlisting}

Job states you will see include \texttt{COMPLETED}, \texttt{FAILED}, \texttt{TIMEOUT} (it hit \texttt{-\/-time}), and \texttt{OUT\_OF\_MEMORY}.

\section{Start small}\label{start-small}

Test on a tiny case before running at scale. A good first job proves that it runs on a compute node (\texttt{hostname}), has the cores you asked for (\texttt{nproc}), finds your software, and can see your files. NI-HPC\textquotesingle s \href{https://ni-hpc.github.io/nihpc-documentation/Green\%20HPC/}{Green HPC page}\footnote{\texttt{https://ni-hpc.github.io/nihpc-documentation/Green\%20HPC/}} makes the same point on carbon grounds: failed jobs waste energy too.

